\section{Proposal B: AND-Tree Optimization Search}
\label{andopt:section}

AND-tree search divides the optimization problem into
assignment cases and compares their results. Unlike the
OR-tree stage, finding one valid timetable is not enough:
every case must be explored or shown unable to improve
the best timetable found. Branch-and-bound allows us
to skip such cases, although the worst-case search
size remains exponential.

\subsection{Search model}
\label{sec:and-model}

The model is $A_{\land}=(S_{\land},T_{\land})$.

\paragraph{Problem descriptions.}
Each problem description is a tuple
$(\mathit{assign},U)$, where $\mathit{assign}$ is a
partial timetable and $U$ contains its unassigned
components:
\[
\mathit{Prob}_{\land}
=\{(\mathit{assign},U):
  \mathit{assign}\text{ is a partial assignment},\
  U=X\setminus\operatorname{dom}(\mathit{assign})\}.
\]
The task at a node is to examine the timetables that extend
its partial assignment and determine whether they improve
the best timetable found so far.

\paragraph{Division relation.}
For a consistent unfinished assignment, choose
$x\in U$ and let
$D_{\mathit{assign}}(x)=\{s_1,\ldots,s_r\}$
be its currently legal slots.
When $r>0$, define
\[
\mathit{Div}(pr,pr_1,\ldots,pr_r),
\]
where
\[
\begin{aligned}
pr&=(\mathit{assign},U),\\
pr_i&=(\mathit{assign}\cup\{x\mapsto s_i\},
       U\setminus\{x\}).
\end{aligned}
\]
Each child assigns the chosen component to a different legal
slot. For example, if a lecture can use either of two slots,
$\mathit{Div}$ creates two cases: one for each slot.
We find the best timetable across these cases. A case can
be skipped only when it has no valid solution or cannot
improve the best timetable already found.

\paragraph{States and extension rules.}
A state is an AND-tree, so
$S_{\land}\subseteq\mathit{Atree}$.
A leaf has the form $(pr,sol)$ and an internal node
has the form $(pr,sol,b_1,\ldots,b_r)$, where
$sol\in\{?,\mathsf{yes}\}$.

An unresolved leaf is expanded using $\mathit{Div}$:
\[
(pr,?)
\longrightarrow
(pr,?,(pr_1,?),\ldots,(pr_r,?)).
\]
When no further search is needed in a case, its leaf is marked resolved:
\[
(pr,?)\longrightarrow(pr,\mathsf{yes}).
\]
These operations can also be applied within a subtree,
leaving the other subtrees unchanged.

Here, $\mathsf{yes}$ means that the comparison task for a
case is complete. This happens when a complete timetable has
been evaluated, no valid completion exists, or the lower bound
shows that the case cannot improve the best timetable.
It does not mean that every case contains a valid timetable.

We use $\mathit{Erw}_{\land}$ for leaf expansion, marking a
completed case, and applying either operation within a subtree.
The transitions used by our process are
\[
T_{\land}
=
\{(q,q')\in S_{\land}\times S_{\land}:
  \mathit{Erw}_{\land}(q,q')\}.
\]

Our control does not use the backtracking operation
$\mathit{Erw}_{\land}^{*}$. Each division includes every
legal slot for the chosen component, so no alternative
division needs to be tried. Completed cases remain marked
$\mathsf{yes}$ while the control selects another unresolved
leaf.

\subsection{Search process}
\label{sec:and-process}

The process is
$P_{\land}=(A_{\land},\mathit{Env}_{\land},K_{\land})$.
The environment contains the scheduling instance,
the best valid timetable found, its score, and fixed
component and slot orders for breaking ties.

\paragraph{Best timetable.}
The OR-tree result initializes
\[
\mathit{assign}_{\mathrm{best}}=\mathit{assign}_0,
\qquad
B=\mathrm{Eval}(\mathit{assign}_{\mathrm{best}}).
\]
Whenever a complete valid timetable has a score below
$B$, it replaces $\mathit{assign}_{\mathrm{best}}$
and updates $B$.

The initial timetable is a benchmark, not a set of
fixed choices. The optimization search starts with
only the mandatory assignments.

\paragraph{Lower bound.}
Within a branch, assigned components keep their slots.
Their missed-preference penalties therefore cannot be removed
by assigning the remaining components. We use their sum as
a lower bound (LB):
\[
\mathrm{LB}(\mathit{assign})
=
\sum_{x\in\operatorname{dom}(\mathit{assign})}
\ \sum_{\substack{s\in\mathit{Slots}\\
s\neq\mathit{assign}(x)}}
\mathit{preference}(x,s).
\]
Every complete timetable in this branch has at least this
penalty, since all other penalties are nonnegative.
Complete timetables are evaluated using the full
$\mathrm{Eval}$ function.

\paragraph{Pruning.}
If $\mathrm{LB}(\mathit{assign}) \ge B$, the branch cannot
improve the best timetable and is marked resolved.
Equality can also be pruned because we need only one
optimal timetable.

\paragraph{Search control.}
The function $f_{\mathrm{leaf}}$ selects the deepest
unresolved leaf, with ties broken by child order.
For that leaf, $f_{\mathrm{trans}}$ selects the
following operation:
\begin{enumerate}
  \item If the partial assignment is inconsistent,
  or an unassigned component has no legal slot,
  resolve the case as infeasible.

  \item If the assignment is complete and valid,
  evaluate it, update the best timetable if its
  score is lower, and resolve the case.

  \item Otherwise, if
  $\mathrm{LB}(\mathit{assign})\ge B$,
  resolve the case by pruning.

  \item Otherwise, select an unassigned component
  using minimum remaining values and apply
  $\mathit{Div}$ to generate all its legal slot cases.
\end{enumerate}
MRV ties use the fixed component order. Children
are ordered by increasing lower bound, with ties
broken by the fixed slot order.

\subsection{Search instance and stopping condition}
\label{sec:and-instance}

Let $\mathit{assign}_{\mathrm{fix}}$ contain the
mandatory assignments, including triggered special
bookings. The search instance is
$\mathit{Inst}_{\land}=(q_0,G_{\land})$, where
\[
q_0=
((\mathit{assign}_{\mathrm{fix}},
  X\setminus\operatorname{dom}(
    \mathit{assign}_{\mathrm{fix}})),?).
\]
The search finishes when every leaf is resolved:
\[
G_{\land}(q)=\mathsf{yes}
\iff
\forall (pr,\mathit{sol})\in\operatorname{Leaves}(q):
\mathit{sol}=\mathsf{yes}.
\]
The search has finitely many assignment cases, and resolved
cases are not explored again. If allowed to finish -- time limit not exceeded -- every
case has either been checked or shown unable to improve
the best timetable. The returned timetable is therefore
optimal.

The process returns
$\mathit{assign}_B=\mathit{assign}_{\mathrm{best}}$.

\subsection{Worked AND-tree derivation}
\label{andopt:demonstration}

Consider three lectures $x,y,z$ of different,
non-500-level courses. Each can use either
$s_1=$ MWF 9:00--10:00 or
$s_2=$ MWF 10:00--11:00.
Both slots have capacity two and minimum zero,
and $x$ and $y$ must not overlap.
There are no mandatory assignments or other hard
restrictions affecting these lectures.

The only nonzero preferences are
\[
\mathit{preference}(x,s_1)=4,\qquad
\mathit{preference}(y,s_1)=2,\qquad
\mathit{preference}(z,s_1)=1.
\]
All other soft penalties are zero.
The initial valid timetable is
\[
\mathit{assign}_0
=\{x\mapsto s_2,y\mapsto s_1,z\mapsto s_2\},
\qquad B=5.
\]

At the root, MRV selects $x$ using the component
order $x<y<z$. We apply
\[
\mathit{Div}(pr_0,pr_1,pr_2),
\]
where
\[
\begin{aligned}
pr_0&=(\varnothing,\{x,y,z\}),\\
pr_1&=(\{x\mapsto s_1\},\{y,z\}),\\
pr_2&=(\{x\mapsto s_2\},\{y,z\}).
\end{aligned}
\]
The tree changes from
\[
(pr_0,?)
\longrightarrow
(pr_0,?,(pr_1,?),(pr_2,?)).
\]
The child bounds are 0 and 4, so the control explores
$pr_1$ first.

Since $x$ occupies $s_1$, only $s_2$ is legal for $y$.
The division $\mathit{Div}(pr_1,pr_3)$ creates
\[
pr_3=(\{x\mapsto s_1,y\mapsto s_2\},\{z\}),
\]
whose lower bound is 2.

Both slots remain legal for $z$.
Applying $\mathit{Div}(pr_3,pr_4,pr_5)$ creates
the following complete cases:
\[
\begin{array}{c|ccc|c}
\text{Case} & x & y & z & \mathrm{Eval}\\
\hline
pr_4 & s_1 & s_2 & s_1 & 2\\
pr_5 & s_1 & s_2 & s_2 & 3
\end{array}
\]
The control evaluates $pr_4$ first and updates $B$ from
$5$ to $2$. Case $pr_5$ is valid but does not improve this
score. Both leaves are marked $\mathsf{yes}$.

The remaining root case $pr_2$ has lower bound $4$.
Since $4 \ge B=2$, it is pruned and marked $\mathsf{yes}$.
The leaves $pr_4$, $pr_5$, and $pr_2$ are now all resolved,
so $G_{\land}(q)=\mathsf{yes}$.

The process returns
\[
\mathit{assign}_B
=\{x\mapsto s_1,y\mapsto s_2,z\mapsto s_1\},
\qquad
\mathrm{Eval}(\mathit{assign}_B)=2.
\]
The unexplored descendants of $pr_2$ cannot improve
this score, so the result is globally optimal.